% Phase 2.5b manuscript, ADS-Cascade / S2AND mechanism-selection study.
% Venue-neutral, generic article class, no journal style file.
% Drafting order (phase2/reports/OUTLINE.md): 4, 5, 6, then 7, then 3, 8, 9,
% then 1, 2, then the abstract.
\documentclass[11pt]{article}

\usepackage[T1]{fontenc}
\usepackage[utf8]{inputenc}
\usepackage[margin=1in]{geometry}
\usepackage{amsmath}
\usepackage{booktabs}
\usepackage{graphicx}
\usepackage{natbib}
\usepackage[hidelinks]{hyperref}

\graphicspath{{../reports/figs/}}
\bibliographystyle{plainnat}

\setlength{\parskip}{0.25\baselineskip}

\title{Prediction is not selection: a pre-registered study of design-time
mechanism choice in author name disambiguation}

\author{Maitreya Sapariya\thanks{ORCID 0009-0003-9346-3775.}}

\date{}

\begin{document}
\maketitle

% ---------------------------------------------------------------------------
% Abstract: drafted last, per OUTLINE.md.
% ---------------------------------------------------------------------------

% ---------------------------------------------------------------------------
% 1. Introduction: drafted in the fifth block.
% ---------------------------------------------------------------------------

% ---------------------------------------------------------------------------
% 2. Related work: drafted in the fifth block.
% ---------------------------------------------------------------------------

% ---------------------------------------------------------------------------
% 3. Problem setting and constructs: drafted in the fourth block.
% ---------------------------------------------------------------------------

\section{A comparison that cannot have a winner}
\label{sec:degenerate}

\subsection{The containment argument}
\label{sec:containment}

Consider the comparison a practitioner reaches for first. Two mechanisms
resolve a held-out name mention against the mentions already recorded under
the same name block. A precision-first \emph{rule} looks for a training
record whose surface form is identical to the test form, assigns the author
carried by that form, and abstains when no such record exists. A
coverage-first \emph{retriever} scores every distinct training form against
the test form with a string similarity function, takes the nearest one, and
assigns its author. Rules against retrieval, exact against fuzzy. Which one
should you build?

Suppose the test form occurs in the block's history. A normalized similarity
function returns its maximum, 1.0, when two strings are identical, and no
other training form can score above the maximum. A retriever that takes the
single nearest form therefore takes the test form itself, and if the two
mechanisms share a tie-break rule over candidates carrying the same form,
they resolve the same author. Whenever the rule fires, the retriever returns
what the rule returns. When the rule abstains, the retriever answers anyway.

The rule's set of correct predictions is a subset of the retriever's. The
rule cannot be right where the retriever is wrong. Under any winner
definition based on accuracy, the verdict ``the rule wins'' is not merely
rare. It is unreachable.

The argument does not depend on the similarity function, since every function
in the normalized string similarity family is maximized by identity, nor on
the acceptance cutoff, since a cutoff cannot exceed 1.0 and so cannot exclude
the identical form. What it does depend on is small and easy to miss: the
retriever takes the single nearest form, and both mechanisms draw their
candidates from the same history.

\subsection{What we measured}
\label{sec:degeneracy-measured}

We did not derive the containment first and then confirm it. We
pre-registered the hypotheses, the mechanisms, the bands, the margins, the
splits and the falsification table; we froze the calibration constants in a
separate commit; we ran the whole analysis; and we read the degeneracy off
the output. That ordering is why this section is a result and not a remark.

Across all 16{,}136 evaluation blocks and 199{,}713 held-out signatures, the
number of signatures on which exact match is correct and fuzzy retrieval is
incorrect is \textbf{zero}. Zero in the Levenshtein arm and zero in the
Jaro-Winkler arm. The cutoff had been calibrated once on a disjoint split,
landing at 0.50 in both metric families; no value on the pre-registered grid
from 0.50 to 0.95 changes the count, because none of them can.

Two pre-registered diagnostics collapsed along with the winner. The pooled
accuracy gap, $-0.0196$ over the primary analysis population, reduces exactly
to the coverage difference, because every signature exact match gets right
fuzzy retrieval also gets right: exact coverage 0.9664 against fuzzy coverage
0.9985 on the evaluation split. The mechanism agreement rate, pre-registered
precisely to tell genuine practical equivalence apart from structural
coincidence, came out at 0.9679, which is exact coverage plus the fuzzy
abstention rate, $0.9664 + 0.0015$, to every printed digit. The diagnostic
written to detect the structural case turned out to compute it.

We report the order of discovery because it is the part that carries
information. The degeneracy is easy to follow once someone states it in a
paragraph. It was, in our hands, hard enough to see in advance that it
survived a pre-registration careful enough to fix a deterministic tie-break
chain for both mechanisms, and it was found by running the full frozen
analysis and looking at an empty cell.

\subsection{The general form}
\label{sec:general-form}

The containment argument is about one pair. The trap is wider. Take
\emph{any} pair of mechanisms that assigns an author drawn from the block's
own history, whatever the machinery inside them. Two regions of the design
space are empty before any data is collected.

\textbf{Where the history names one author, a precision-first rule cannot
win.} If every training record in the block belongs to one author, then any
mechanism that returns an author from a non-empty candidate pool returns that
author. The rule returns the same author whenever it fires. Its only
available difference is to abstain, and abstention scores as incorrect. So
the rule can lose there and can tie there, and cannot win.

\textbf{Where the history is written one way, no mechanism can win.} If the
block's training records carry a single distinct surface form, every
mechanism sees the same form and the same evidence attached to it. They all
return the same author. The comparison is a tie by construction.

Both statements are visible in our calibration-split probe of the redesigned
pair described in Section~\ref{sec:repaired-pair}. Counting held-out
signatures on which the rule is right and the retriever wrong, across the
five consistency bands crossed with the four surface-variability bands, the
count is exactly zero in all four cells of the top consistency band and
exactly zero in all five cells of the zero-variability column. The other two
candidate pairs we probed show the same two lines of zeros. A contest is
possible only where the history is both impure and surface-varied, which on
the calibration split is 41\% of held-out signatures.

This is a trap practitioners fall into, not a theorem anyone needs. Anyone
comparing a string-matching baseline against a retrieval or clustering method
over the same recorded history is inside this family, and per-instance
algorithm selection over such a portfolio inherits the problem whole: two
regions of the selection space hold nothing to select. The check is cheap.
Before running the comparison, count the items on which mechanism A is
correct and mechanism B is not. If that count is zero, or zero throughout an
entire stratum, the design produced it and the data had no say. Report those
regions in the design section, where they are a property of the comparison,
rather than in the results, where they read as a defect.

\section{A pair that can}
\label{sec:repaired-pair}

\subsection{Three candidates, and why a k-nearest vote}
\label{sec:candidates}

The repair cannot be a different similarity function, because the containment
holds for all of them. Mechanism B has to stop being a single-nearest
retriever. We probed three ways of doing that on the calibration split, the
15\% of blocks the pre-registration reserves for design choices, covering
2{,}450 blocks and 35{,}732 held-out signatures. Mechanism A, the
exact-evidence rule, is unchanged in all three pairs.

\emph{Radius vote.} Pool every training record whose form is within
similarity $r$ of the test form, and take the most frequent author in the
pool. \emph{Normalized-key vote.} Pool every training record whose form
carries the same ASCII-folded, lowercased, letters-only key as the test form.
\emph{k-nearest vote.} Pool the nearest training records by form similarity,
whole forms at a time, until at least $k$ records are pooled, and vote.

All three are two-sided in the probe. We adopted the k-nearest vote at
$k = 5$, for three reasons in order of weight. It gives the most balanced
contest: 546 signatures where the rule is right and the vote wrong against
744 the other way, with pooled accuracies 0.8234 and 0.8290. Rank beats
radius on this substrate, because every name in a block shares a surname, so
absolute similarity has no natural scale and a fixed radius means different
things in different blocks; the radius family confirms it, with the rule-side
count running from 1{,}461 at $r = 0.70$ down to 0 at $r = 0.95$, so the
whole finding would have been a function of $r$. And the k-nearest vote never
abstains, which puts abstention entirely on the rule's side and separates the
two arms along exactly the precision-against-coverage axis the paper is
about.

\subsection{The calibration trap}
\label{sec:calibration-trap}

Read the accuracy column of the same probe and a second problem appears. The
k-nearest vote's pooled accuracy falls monotonically as $k$ grows: 0.8390,
0.8354, 0.8290 and 0.8188 at $k$ of 2, 3, 5 and 10. The radius vote's
accuracy rises monotonically as $r$ grows, reaching 0.8239 at $r = 0.95$,
where its rule-side count is 0. In both families the accuracy-maximizing
parameter is the one nearest the single-nearest limit, and the single-nearest
limit is the degeneracy of Section~\ref{sec:degenerate}.

Our original pre-registration calibrated the retriever's one parameter by
picking the value with the highest pooled accuracy on the calibration split.
That protocol, applied unchanged to either parameterized pair, selects the
degenerate parameter and walks the design straight back into a comparison
that cannot have a winner. We therefore voided it. $k$ is fixed by definition
at 5, disclosed as a calibration-split choice made for contest balance rather
than for accuracy, with a mandatory sensitivity curve over $k$ in
$\{2, 3, 5, 10\}$ reported beside every outcome.

The general warning is worth more than the fix. When one member of a compared
pair contains the other in a limit, tuning that member for accuracy is tuning
it toward containment. A protocol that selects a parameter by maximizing the
tuned arm's own performance will find that limit, and will then report a
degenerate comparison as a well-calibrated one.

\section{Experimental design and pre-registration}
\label{sec:design}

\subsection{Data, blocking and splits}
\label{sec:data}

The substrate is S2AND \citep{s2and2021}, which unifies eight prior author
name disambiguation datasets into one format with a reference implementation.
A \emph{signature} is one author mention on one paper, carrying a name string
and context fields. Signatures are grouped into \emph{blocks} by a normalized
name key, frozen before any measurement, and ground truth assigns each
signature to a real author. We use seven of the eight constituent datasets:
aminer, arnetminer, inspire, kisti, pubmed, qian and zbmath. Medline is
excluded because the release ships no clusters file for it.

Table~\ref{tab:population} accounts for every signature from the corpus down
to the analysis population. Blocks with a single signature have no history to
match against, so they sit outside the live set. Live blocks are split at
block level with seed 20260918 into a 15\% calibration split and an 85\%
evaluation split. The calibration split is used once, for the band edges and
for the cutoff of the mechanism we later retired, and never for a reported
hypothesis number. The evaluation split produces every result below. Inside
each evaluation block, signatures are assigned to a training half and a test
half with the same seed, by
$n_{\mathrm{train}} = \max(1, \min(n - 1, \lfloor 0.7n \rfloor))$. The
training signatures are the block's historical knowledge base. The test
signatures are scored, and no test signature is visible to any signal, band
edge, cutoff or mechanism.

The primary analysis population is evaluation blocks with at least two
training signatures and at least one test signature, that is blocks of three
or more signatures. Below that, consistency is 1.0 by construction on a
one-signature history and surface variability is undefined, so neither signal
carries information. Blocks of exactly two are a named and counted exclusion,
reported as their own stratum and used to decide nothing.

\begin{table}[htbp]
\centering
\caption{Population accounting, from the corpus to the primary analysis
population. Both exclusions are named and counted rather than dropped.
Source: \texttt{reports/phase2e\_boundary.md}, generated by
\texttt{phase2e\_run.py}.}
\label{tab:population}
\begin{tabular}{lrrr}
\toprule
Stratum & Blocks & Signatures & Share of signatures \\
\midrule
All signatures & 26{,}939 & 766{,}133 & 100.00\% \\
Singleton blocks, outside the live set & 7{,}955 & 7{,}955 & 1.04\% \\
Live blocks (2 or more signatures) & 18{,}984 & 758{,}178 & 98.96\% \\
\quad Calibration split (15\%) & 2{,}848 & 116{,}308 & 15.18\% \\
\quad Evaluation split (85\%) & 16{,}136 & 641{,}870 & 83.78\% \\
\quad\quad Primary analysis population & 13{,}902 & 637{,}402 & 83.20\% \\
\quad\quad Block-size-2 stratum, named exclusion & 2{,}234 & 4{,}468 & 0.58\% \\
\bottomrule
\end{tabular}
\end{table}

Every reported number comes from the 13{,}902 blocks and 197{,}479 test
signatures of the primary analysis population. A cell, band or dataset
holding fewer than 30 blocks is printed with its count, marked under-powered,
interpreted nowhere and used to decide nothing. A dataset qualifies for the
generalization rule if it contributes at least 100 blocks.

\subsection{The two signals and their bands}
\label{sec:signals}

Both signals are computed from the training half of each block only. That is
what design-time means here, and it is load-bearing: a signal that needs the
test items, or their labels, is not a signal a practitioner can act on before
building anything.

\textbf{ADS}, the consistency signal, is the share of a block's training
signatures belonging to its single most frequent author. It runs on
$(0, 1]$ and equals 1.0 when the history names one person. High ADS means low
homonymity. Its bands are fixed by the definition of purity rather than read
off a distribution: $[0, 0.5)$, $[0.5, 0.7)$, $[0.7, 0.9)$, $[0.9, 1.0)$ and
exactly 1.0.

\textbf{SV}, the surface-variability signal, is the expected normalized
Levenshtein distance between two training signatures of the block drawn at
random without replacement, computed on the frozen surface form and weighted
by signature. Higher SV means less surface stability. Its scale is bounded
well below 1 by construction, because every name in a block shares a surname,
so fixed bands on the nominal scale would be meaningless. SV = 0, a history
written one way, is structural and is its own band. The other three bands are
the terciles of the strictly positive values, with cut points taken from the
calibration split at 0.102644 and 0.195066 and committed before any
evaluation number existed.

Two controls travel with SV throughout, because the largest single source of
within-block name variation in bibliographic data is forename initialization,
which is the variable \citet{kimkim2020} manipulate. \textbf{SV\_resid} is
the same function applied to an initialization-canonical form, in which the
first name and each middle token are reduced to their first letter, so that
``John Smith'' and ``J. Smith'' become the identical string and
initialization contributes exactly zero. What remains is spelling,
diacritics, hyphenation, transliteration, middle-name presence and suffixes.
Its band edges, calibrated the same way, are 0.036782 and 0.082418.
\textbf{INIT} is the share of a block's training signatures whose given-name
field is an initial only, banded at INIT = 0, $0 < \mathrm{INIT} < 0.5$ and
$\mathrm{INIT} \geq 0.5$, and the SV analysis is repeated inside each INIT
band.

The two axes were measured alive and separable before any mechanism was
built. ADS and SV correlate at Spearman $-0.2612$ across the 18{,}984 live
blocks, and each of our measures tracks its S2AND counterpart strongly while
tracking the other weakly: Spearman $-0.818$ for ADS against S2AND's
homonymity facet, $+0.743$ for SV against its synonymity facet.

\subsection{Mechanisms, metric and the winner}
\label{sec:winner-construct}

\textbf{Mechanism A, the exact-evidence rule.} If the test signature's
surface form occurs among the training forms, assign the most frequent author
among training signatures carrying that exact form. Otherwise abstain. No
parameter.

\textbf{Mechanism B, the k-nearest vote.} Pool the nearest training
signatures by similarity of their surface form to the test form, whole forms
at a time, until the pool holds at least $k = 5$ training signatures; if the
block holds fewer, the pool is the whole block. Vote the pooled signatures'
authors, weighted by signature, and assign the winner. B does not abstain.

Both mechanisms are untrained lookups over the block's own history, and both
use the same frozen tie-break chain: most frequent author inside the pooled
evidence, then most frequent author inside the block, then the
lexicographically smallest author identifier. Sharing the chain matters,
because giving the two arms different tie-breaks would turn part of the
winner contest into a comparison of tie-break rules.

\textbf{Metric.} Whole-set accuracy over the block's test signatures, with
abstention counted as incorrect. Coverage, the share of test signatures not
abstained on, is reported separately for each mechanism wherever accuracy is
reported. The asymmetry that the rule abstains and the vote does not is a
result we keep visible rather than normalize away.

\textbf{Winner.} Per block, we take a paired bootstrap on the rule's accuracy
minus the vote's accuracy, resampling the block's test signatures with
replacement, $B = 2000$ draws, with both mechanisms scored on the same
resample, and read the 95\% percentile interval. Per band and per grid cell
we pool the band's test signatures and bootstrap by resampling \emph{blocks}
with replacement, because the block is the independent unit. Four verdicts,
mutually exclusive, against a practical-equivalence margin $\delta$:
\emph{rule wins} if the interval's lower bound exceeds $+\delta$,
\emph{k-NN wins} if its upper bound falls below $-\delta$, \emph{tie} if the
interval lies entirely inside $[-\delta, +\delta]$, and \emph{undecided}
otherwise.

Tie and undecided are kept apart throughout, and we state the distinction
here, before any result, because merging them is how a low-powered comparison
comes to read as a positive finding of equivalence. A tie means equivalence
was demonstrated. Undecided means the evidence did not reach a decision. We
also report the verdict distribution against test-set size, so a reader can
see how much of any no-winner region is one and how much is the other.

The primary margin is $\delta = 0.05$ accuracy points, fixed in advance as a
difference a practitioner would act on, with $\delta = 0.02$ and
$\delta = 0.03$ pre-registered as a robustness curve. An outcome that holds
at one margin and flips at another is reported as margin-dependent and
treated as Partial, never Supported. The same rule applies across the $k$
curve. And Inconclusive is evaluated first: if a hypothesis's pre-registered
power floor is met, the outcome is Inconclusive regardless of what the other
rows would have said. A set of bands that are all tie or all undecided is
never read as support.

\subsection{Pre-registration and its amendment}
\label{sec:prereg-integrity}

The text of this subsection was written during the amendment pass, before the
run it defends was executed, and locked in the project plan at that point. It
is reproduced here as written.

\begin{quote}
We pre-registered the hypotheses, the winner definition, the analysis bands,
and the falsification table before computing any evaluation metric, and we
froze the calibration constants in a separate commit before any evaluation
number existed. The pre-registration and every subsequent amendment are in
the public repository with full history.

We amended the pre-registration once, after two exploratory runs, and we
state that plainly rather than presenting the amended version as the
original. The amendment was forced by two structural facts about the
mechanisms, not by an unwelcome effect size. First, when a test name is
present in a block's history it sits at the maximum of any similarity
function, so a single-nearest retriever returns whatever exact match returns;
exact match is therefore a strict sub-case of single-nearest retrieval, and
the two cannot contest a winner. Second, a block whose history names one
author returns that author from any non-empty candidate pool, so at maximum
consistency the only move a mechanism can make is to abstain. The first fact
required replacing the second mechanism with a k-nearest vote that does not
reduce to exact match. The second required scoping the consistency hypothesis
to whether the choice is decidable at all, rather than which mechanism wins
at every consistency level.

Three properties keep the amendment honest. The original hypotheses and the
conditions of the runs that motivated the change remain in the
pre-registration, unedited and dated. No accuracy figure for the k-nearest
mechanism existed at the time the amended hypotheses were written; the
amendment predates the run it governs. And the one component of the amended
consistency hypothesis that is guaranteed by construction, the decline in
decided winners as consistency approaches its maximum, is labeled as
structural wherever it appears; the empirical content we test lies below that
maximum, in the region where the outcome is not fixed in advance.
\end{quote}

Three points of record, checked at writing against the committed repository
rather than against memory.

\textbf{What ``amended once'' means.} The repository holds two dated appends
to the pre-registration, not one. The first is the amendment itself, which
replaced the second mechanism and rewrote the two selection hypotheses. The
second is a short addendum that replaced one clause's measure of which signal
drives decidability, because the original measure was an order statistic
taken over a grid with five rows and four columns, and so favoured the axis
with more bands, which was the axis the hypothesis needed to win. Both were
committed in the same pass, within half an hour of each other, and both
precede the run they govern. They are one amendment event recorded in two
documents. The addendum tightens a rule and relaxes none.

\textbf{What the two runs before the amendment were.} The first was not
exploratory. It was the pre-registered run, executing the frozen
pre-registration exactly, and it is what measured the degeneracy in
Section~\ref{sec:degeneracy-measured}. The second was exploratory by
construction: a design probe confined to the calibration split, which
produced no evaluation-split number and evaluated no hypothesis.

\textbf{The ordering claim, against commit timestamps.} The pre-registration
was committed before the run that executed it. The calibration constants,
including the band edges reused unchanged in the final run, were committed in
their own earlier commit, before any evaluation-split accuracy existed. The
amendment was committed just under two hours before the run it governs and
the addendum about eighty minutes before it, and no prediction of the
k-nearest vote on an evaluation block existed anywhere in the repository at
either moment.\footnote{Commits, in order:
\texttt{8869b82} freezes the pre-registration; \texttt{5e13e38} freezes the
calibration constants; \texttt{3d86959} is the first run; \texttt{c20239a} is
the calibration-split design probe; \texttt{5f452ad} and \texttt{3e79a48} are
the amendment and its addendum; \texttt{9e3d921} is the run reported here.}
The first evaluation-split number for the amended mechanism appears in the
run commit itself. That run opens with a regression check, reported in
Section~\ref{sec:two-sided}, which reproduces calibration-split counts frozen
before the amendment was written and halts the run if they do not match.

% ---------------------------------------------------------------------------
% 7. Results: drafted in the third block.
% ---------------------------------------------------------------------------

% ---------------------------------------------------------------------------
% 8. Discussion, 9. Limitations, 10. Conclusion: drafted in the fourth block.
% ---------------------------------------------------------------------------

\bibliography{references}

\end{document}
